% Options for packages loaded elsewhere
\PassOptionsToPackage{unicode}{hyperref}
\PassOptionsToPackage{hyphens}{url}
\documentclass[
]{article}
\usepackage{xcolor}
\usepackage[margin=1in]{geometry}
\usepackage{amsmath,amssymb}
\setcounter{secnumdepth}{-\maxdimen} % remove section numbering
\usepackage{iftex}
\ifPDFTeX
  \usepackage[T1]{fontenc}
  \usepackage[utf8]{inputenc}
  \usepackage{textcomp} % provide euro and other symbols
\else % if luatex or xetex
  \usepackage{unicode-math} % this also loads fontspec
  \defaultfontfeatures{Scale=MatchLowercase}
  \defaultfontfeatures[\rmfamily]{Ligatures=TeX,Scale=1}
\fi
\usepackage{lmodern}
\ifPDFTeX\else
  % xetex/luatex font selection
\fi
% Use upquote if available, for straight quotes in verbatim environments
\IfFileExists{upquote.sty}{\usepackage{upquote}}{}
\IfFileExists{microtype.sty}{% use microtype if available
  \usepackage[]{microtype}
  \UseMicrotypeSet[protrusion]{basicmath} % disable protrusion for tt fonts
}{}
\makeatletter
\@ifundefined{KOMAClassName}{% if non-KOMA class
  \IfFileExists{parskip.sty}{%
    \usepackage{parskip}
  }{% else
    \setlength{\parindent}{0pt}
    \setlength{\parskip}{6pt plus 2pt minus 1pt}}
}{% if KOMA class
  \KOMAoptions{parskip=half}}
\makeatother
\usepackage{graphicx}
\makeatletter
\newsavebox\pandoc@box
\newcommand*\pandocbounded[1]{% scales image to fit in text height/width
  \sbox\pandoc@box{#1}%
  \Gscale@div\@tempa{\textheight}{\dimexpr\ht\pandoc@box+\dp\pandoc@box\relax}%
  \Gscale@div\@tempb{\linewidth}{\wd\pandoc@box}%
  \ifdim\@tempb\p@<\@tempa\p@\let\@tempa\@tempb\fi% select the smaller of both
  \ifdim\@tempa\p@<\p@\scalebox{\@tempa}{\usebox\pandoc@box}%
  \else\usebox{\pandoc@box}%
  \fi%
}
% Set default figure placement to htbp
\def\fps@figure{htbp}
\makeatother
\setlength{\emergencystretch}{3em} % prevent overfull lines
\providecommand{\tightlist}{%
  \setlength{\itemsep}{0pt}\setlength{\parskip}{0pt}}
\usepackage{bookmark}
\IfFileExists{xurl.sty}{\usepackage{xurl}}{} % add URL line breaks if available
\urlstyle{same}
\hypersetup{
  pdftitle={Acceso a la Justicia y Litigio Constitucional},
  hidelinks,
  pdfcreator={LaTeX via pandoc}}

\title{Acceso a la Justicia y Litigio Constitucional}
\author{}
\date{\vspace{-2.5em}}

\begin{document}
\maketitle

Demián Zayat
\href{mailto:dzayat@derecho.uba.ar}{\nolinkurl{dzayat@derecho.uba.ar}}
Ayudante Tamara London
\href{mailto:tamaralondon@derecho.uba.ar}{\nolinkurl{tamaralondon@derecho.uba.ar}}

\subsection{Programa}\label{programa}

En este curso analizaremos el derecho de acceder a la justicia, su
fuente constitucional, sus garantías de funcionamiento y sus mayores
problemas y dificultades para su ejercicio.

Nuestro sistema jurídico presupone que cualquier persona podrá realizar
una demanda ante los tribunales cuando sus derechos hayan sido
afectados, vulnerados o amenazados. Este presupuesto asume que el Poder
Judicial podrá entender y propender a la solución de cualquier conflicto
de derechos que se le presente. Sin embargo, existen diversas razones
que hacen que los tribunales de justicia y la protección judicial no sea
accesible de modo igualitario para todas las personas. Algunas tendrán
muchas menos dificultades para hacerlo y otras deberán vencer muchas más
barreras para poder hacer llegar su reclamo a los jueces.

Existen aún hoy algunas leyes y normativas que hacen que este acceso sea
especialmente dificil para ciertos grupos. En esos casos, la solución
podrá ser --paradójicamente-- vía un proceso constitucional.

Sin embargo, también existen supuestos de grupos con mayores
dificultades para hacer valer sus derechos ya no por una normativa
expresa sino por una situación social, y diversas prácticas de los
operadores judiciales y estatales en general. Estos casos son más
dificiles de identificar, aunque inciden de un modo de mayor
performatividad, siendo complicado desestabilizar la situación
existente. Muchas veces puede verse que la respuesta judicial resulta
patriarcal, xenófoba, racista, hegemónica, etc.

En este curso haremos una revisión sobre la normativa constitucional
involucrada en el derecho de acceso a la justicia, tomaremos nota de las
garantías procesales que establece la constitución para su eficacia, y
analizaremos de modo crítico y con casos concretos cómo y por qué
existen grupos identitarios de población que suele tener mayores
dificulades de acceso a los tribunales.

\subsection{Objetivo}\label{objetivo}

El objetivo del curso es visibilizar, discutir y analizar de modo
crítico las diversas barreras que afectan al acceso a la justicia, sus
aspectos centrales, y su dimensión estructural. La idea es poder
analizar problemas jurídicos desde el punto de vista del acceso a la
justicia y decidir estrategias para superarlas.

\subsection{Metodología}\label{metodologuxeda}

Será un curso con un componente teorico sobre el derecho de acceso a la
justicia y sus garantías de funcionamiento, y práctico en el análisis de
la situación de los problemas que los grupos identitarios tienen en
particular. Este aspecto práctico partirá del análisis de diversas
investigaciones empíricas sobre cuáles son estas dificultades, y la
discusión sobre cómo fue resuelta tanto por los tribunales locales como
regionales o internacionales.

Se espera que les estudiantes puedan leer la bibliografía indicada de
modo previo a cada clase, y que los encuentros funcionen no para repetir
sino para debatir sobre la literatura indicada.

Habrá un único examen presencial al final de la cursada.

\subsection{Cronograma y
bibliografía}\label{cronograma-y-bibliografuxeda}

\subsubsection{Semana 01, viernes 13/03. Unidad 1.
Presentación}\label{semana-01-viernes-1303.-unidad-1.-presentaciuxf3n}

Presentación de los objetivos y metodología del curso. Introducción al
derecho de acceso a la justicia. Dimensión instrumental y dimensión
sustantiva del derecho. Movimiento de Acceso a la Justicia.
Reconocimiento constitucional.

\begin{itemize}
\tightlist
\item
  Maurino, Gustavo y Sucunza Matías, ``Acceso a la Justicia'', en
  Gargarella- Guidi \emph{Comentarios de la Constitución de la Nación
  Argentina}, tomo 2, pp 895-929, puntos 1 y 2,
  \href{https://drive.google.com/file/d/1On39ejMciqQNjJIpz1fPseKak_KCjuuF/view?usp=sharing}{aca}
\end{itemize}

\subsubsection{Semana 02, martes 17/03 Unidad 1. Estándares
internacionales}\label{semana-02-martes-1703-unidad-1.-estuxe1ndares-internacionales}

\begin{itemize}
\item
  Maurino, Gustavo y Sucunza Matías, ``Acceso a la Justicia'', en
  Gargarella- Guidi \emph{Comentarios de la Constitución de la Nación
  Argentina}, tomo 2, pp 895-929, puntos 3 y 4,
  \href{https://drive.google.com/file/d/1On39ejMciqQNjJIpz1fPseKak_KCjuuF/view?usp=sharing}{aca}.
\item
  Reglas de Brasilia sobre Acceso a la Justicia de las Personas con
  vulnerabilidad con la actualización de 2018,
  \href{https://brasilia100r.com/wp-content/uploads/2020/07/Reglas-de-Brasilia-actualizaci\%C3\%B3n-2018.pdf}{en
  este link}
\item
  Borrador de Convenio Iberoamericano sobre Acceso a la Justicia
  \href{https://drive.google.com/file/d/1ZvhBoVmS245Mt4_mxBN_GJOdlh1jDEL5/view?usp=sharing}{aca}
\end{itemize}

\subsubsection{Semana 02, viernes 20/03 Unidad 2. Concepto clásico y
estructural}\label{semana-02-viernes-2003-unidad-2.-concepto-cluxe1sico-y-estructural}

Patrocinio letrado. Requerir un patrocinio letrado puede ser difícil en
muchos casos. ¿Es una garantía útil? Ahora, son todes les abogades
iguales? Por qué tener dinero es relevante para el sistema judicial?
Concepciones clásica y estructural del acceso a la justicia.

\begin{itemize}
\item
  Gallanter, ``Por qué los poseedores salen adelante?,
  \href{https://drive.google.com/file/d/1Fj3STLmXxHrJglgGgs5JUtDjXy0DSJ3o/view?usp=sharing}{link}
\item
  Bohmer, ``Igualadores Retóricos'',
  \href{https://drive.google.com/file/d/0B50ljTnhr79kMWQ2OTgzYzMtNGNiNS00YjQxLThjYjYtYzg2YWYyNzc0MTlh/view?usp=sharing&resourcekey=0-QzTl1yGSNgnHMOI_ICHjEg}{aca}
\end{itemize}

\subsubsection{\texorpdfstring{{Semana 03, martes 24/03 Feriado 24 de
marzo. Día de la Memoria, por la verdad y la
justicia}}{Semana 03, martes 24/03 Feriado 24 de marzo. Día de la Memoria, por la verdad y la justicia}}\label{semana-03-martes-2403-feriado-24-de-marzo.-duxeda-de-la-memoria-por-la-verdad-y-la-justicia}

\subsubsection{Semana 03, viernes 27/03. Unidad 3. La Agenda estructural
-
Víctimas}\label{semana-03-viernes-2703.-unidad-3.-la-agenda-estructural---vuxedctimas}

\begin{itemize}
\item
  CELS, ``El acceso a la justicia como una cuestión de derechos
  humanos'', en Informe Anual 2016,
  \href{https://drive.google.com/file/d/1qH6DsnHctmesy6nje2f5NgCyowQ0PkyM/view?usp=sharing}{acá}
  El acceso a la justicia como una cuestión de derechos humanos
\item
  Piqué, María, ``Los derechos de las víctimas de delitos en nuestra
  Constitución'', disponible
  \href{https://drive.google.com/file/d/1NaUqgIjcAosZx6OCPuJzcdG3sbKrqfaC/view?usp=sharing}{aca}
\item
  Ley de víctimas
  \href{https://servicios.infoleg.gob.ar/infolegInternet/anexos/275000-279999/276819/norma.htm}{aca}
\item
  Declaración de Karina Gimenez, a partir del minuto 28 hasta el minuto
  35.
\end{itemize}

\subsubsection{Semana 04, martes 31/03. Unidad 3. La práctica -
Migrantes}\label{semana-04-martes-3103.-unidad-3.-la-pruxe1ctica---migrantes}

En esta unidad comenzaremos a ver cómo son afectados algunos grupos en
particular.

\begin{enumerate}
\def\labelenumi{\alph{enumi}.}
\tightlist
\item
  Migrantes Expulsiones. Caso Vanessa Gomez
\end{enumerate}

\begin{itemize}
\item
  Expediente judicial CAF 067668/2018
  \href{https://drive.google.com/file/d/1MszNbvuQgl56ukSCbql4ggAAKyiia5Bs/view?usp=sharing}{retencion},
  \href{https://drive.google.com/file/d/11RSJ0lXuMKwhQy-gVHfODbc-AH1fBE0j/view?usp=sharing}{captura}
  y
  \href{https://drive.google.com/file/d/1xsXyT2y3F7sGL8yhK8kG6GTfhkYi7nVw/view?usp=sharing}{apelacion}
\item
  Oteiza y Zayat, A treinta años de la CDN,
  \href{https://drive.google.com/file/d/1rZbIfh3834THG1K07moFq1Cgp_zmYtsk/view?usp=sharing}{aca}
\item
  CSJN Fallo ``Barrios Rojas''
  \href{https://drive.google.com/file/d/15DQUBVIVb6feRRGwkksVgjZT9qRLTdOL/view?usp=sharing}{aca}
\item
  CSJN, Fallo ``Zuluaga Celemin'' del 13/10/2022
  \href{https://drive.google.com/file/d/1B0sp1vdW7HV6fYlQxB1nLoCG4GXRYcR8/view?usp=sharing}{aca}
\item
  CSJN Fallo ``Chiemena Uhiara''
  \href{https://drive.google.com/file/d/1m4EdxUp68KgqEPfTHjRH59QR_IDpFAk5/view?usp=sharing}{aca}
\item
  CSJN, Fallo ``Ramírez Balbuena''
  \href{https://drive.google.com/file/d/1Ywkky3GM7QkgfrJzKadus-wxZKdKI1aj/view?usp=sharing}{aca}
  (agregado recién)
\item
  DNU 366/2025
  \href{https://servicios.infoleg.gob.ar/infolegInternet/anexos/410000-414999/413297/norma.htm}{aca}
\end{itemize}

\subsubsection{\texorpdfstring{{Semana 04, viernes 03/04. Feriado -
Viernes
Santo}}{Semana 04, viernes 03/04. Feriado - Viernes Santo}}\label{semana-04-viernes-0304.-feriado---viernes-santo}

\subsubsection{Semana 05, martes 07/04. Unidad 3. La práctica -
Situación de
Calle}\label{semana-05-martes-0704.-unidad-3.-la-pruxe1ctica---situaciuxf3n-de-calle}

\begin{enumerate}
\def\labelenumi{\alph{enumi}.}
\setcounter{enumi}{1}
\tightlist
\item
  Usurpaciones y derecho de defensa
\end{enumerate}

\begin{itemize}
\tightlist
\item
  Ricciardi y Zayat, \emph{El derecho de defensa en los casos de
  usurpaciones en la Ciudad de Buenos Aires: un estudio empírico}, en
  Revista Institucional de la Defensa Pública de la CABA N° 1,
  disponible
  \href{https://drive.google.com/file/d/1KeqFSqkZ2aBrHXkRcIaZCVu6AMSaX7zH/view?usp=sharing}{acá}
\end{itemize}

\subsubsection{Semana 05, viernes 10/04. Unidad 3. La práctica-
Discapacidad}\label{semana-05-viernes-1004.-unidad-3.-la-pruxe1ctica--discapacidad}

\begin{enumerate}
\def\labelenumi{\alph{enumi}.}
\setcounter{enumi}{2}
\tightlist
\item
  Personas con discapacidad
\end{enumerate}

\begin{itemize}
\item
  Corte IDH, caso ``Furlan y familiares c.~Argentina'', disponible
  \href{https://www.corteidh.or.cr/docs/casos/articulos/seriec_246_esp.pdf}{aca},
  y un resumen oficial
  \href{https://www.corteidh.or.cr/docs/casos/articulos/resumen_246_esp.pdf}{aca}
\item
  CSJN, ``Asociación Francesa de Beneficencia''
  \href{https://drive.google.com/file/d/1hVROXu_IFbswD_0UtcBpXPZ0gEbTIiQJ/view?usp=share_link}{aca}
\item
  (opcional) CSJN, ``Institutos Médicos Antártida s/ quiebra s/ inc. de
  verificación (R.A.F. y L.R.H. de F)''
  \href{https://drive.google.com/file/d/1THtTAOVsLhI8y5mQq5mNTLrt-LFdFc-K/view?usp=sharing}{aca}
\item
  CSJN, Ejecuciones contra el Estado
  \href{https://drive.google.com/file/d/1X2oK8aGhqhadJ7YkQCHktaT49Ypk5udZ/view?usp=share_link}{aca}
\end{itemize}

\subsubsection{Semana 06, martes 14/04. Unidad 3. La práctica - Salud
Mental -
Niñeces}\label{semana-06-martes-1404.-unidad-3.-la-pruxe1ctica---salud-mental---niuxf1eces}

\begin{enumerate}
\def\labelenumi{\alph{enumi}.}
\setcounter{enumi}{3}
\tightlist
\item
  Salud mental
\end{enumerate}

\begin{itemize}
\tightlist
\item
  CSJN, ``R. M.J s/Insania'',
  \href{https://drive.google.com/file/d/1c29zw3cQDmwNGuO-GDR15KUq7Xc-UZJb/view?usp=sharing}{aca}
\end{itemize}

\begin{enumerate}
\def\labelenumi{\alph{enumi}.}
\setcounter{enumi}{4}
\tightlist
\item
  Niñeces
\end{enumerate}

\begin{itemize}
\tightlist
\item
  MPT CABA,
  \href{https://drive.google.com/file/d/1lc3H955K8Wx1-7mdtDAJX--CBBcqXfuX/view?usp=sharing}{Violencia
  contra Niños, Niñas y Adolescentes} (solo Capítulos 1 y 2)
\end{itemize}

\begin{enumerate}
\def\labelenumi{\alph{enumi}.}
\setcounter{enumi}{3}
\tightlist
\item
  LGBT Trans
\end{enumerate}

\begin{itemize}
\item
  Informe sombra CEDAW 2016,
  \href{https://drive.google.com/file/d/1Zv0TtTTzgt5N4vSBEo3F7UjiayBQ2BIx/view?usp=sharing}{aca}
\item
  Violencia institucional contra personas trans,
  \href{https://drive.google.com/file/d/1NWDuVTHVpSSWbHTuvP6WPAy6waXvDyXo/view}{ponencia}
\item
  Caso Corte IDH, Vicky Hernandez c/ Honduras
  \href{https://www.corteidh.or.cr/docs/casos/articulos/seriec_422_esp.pdf}{sentencia}
  y
  \href{https://www.corteidh.or.cr/docs/casos/articulos/resumen_422_esp.pdf}{resumen}
  y \href{https://www.youtube.com/watch?v=5AXht8vR6l8}{video} a partir
  del minuto 40.
\item
  Caso de Luz Aime,
  \href{https://www.cosecharoja.org/absolucion-para-luz-le-quieren-dar-perpetua-por-trans-y-pobre/}{aca}
\end{itemize}

\subsubsection{Semana 07, martes 21/04. Unidad 3. La
práctica}\label{semana-07-martes-2104.-unidad-3.-la-pruxe1ctica}

Unidad 3. La práctica - Mujeres

\begin{enumerate}
\def\labelenumi{\alph{enumi}.}
\setcounter{enumi}{4}
\tightlist
\item
  Mujeres
\end{enumerate}

\begin{itemize}
\item
  Pique y Fernandez Valle, ``La garantía de imparcialidad judicial desde
  la perspectiva de género''
  \href{https://drive.google.com/file/d/1atsNG8fNfxcBruDPtPve0N3AmUf6OA3h/view?usp=sharing}{aca}
\item
  Recusación Luz Aime
  \href{https://drive.google.com/file/d/1ZWe1xva3x_Y0ttVg2CZ4eliKBj-8bbBy/preview}{aca}
\item
  CSJN, Fallo Rizzi, 16/05/2024
  \href{https://drive.google.com/file/d/1W6zEuSV0_wshgD8vEc_1HMRiGjlUa-qS/view?usp=sharing}{aca}
\item
  Alegato de cierre Ministerio Público Fiscal
  \href{https://youtu.be/ZKH82JYVLgc}{Luz Aime}
\end{itemize}

\subsubsection{Semana 07, viernes 24/04. Unidad 3. La práctica -
Mujeres}\label{semana-07-viernes-2404.-unidad-3.-la-pruxe1ctica---mujeres}

\begin{enumerate}
\def\labelenumi{\alph{enumi}.}
\setcounter{enumi}{4}
\tightlist
\item
  Mujeres (cont)
\end{enumerate}

\begin{itemize}
\item
  Inecip, casos de violencia de género,
  \href{https://inecip.org/publicaciones/herramientas-jurisprudenciales-para-el-litigio-con-perspectiva-de-genero/}{aca}.
  Leer la introducción y las fichas 1, 2, 3, 5 y 6.
\item
  CEDAW, Observación General 33, disponible
  \href{https://drive.google.com/file/d/1j3suY4Q3JhS5EzbAvs6_K6XSbQsOQ7D7/view?usp=sharing}{acá}
\item
  Ley de Protección integral conta la Mujer
  \href{http://servicios.infoleg.gob.ar/infolegInternet/anexos/150000-154999/152155/texact.htm}{aca}
\end{itemize}

Dejo por aca el video de \href{https://youtu.be/prUqWHKnUd4}{Rita
Segato}

\begin{itemize}
\tightlist
\item
  CEDAW, comunicacion 88/2015 - Sra X c.~Timor Lester
  \href{https://drive.google.com/file/d/1eyLvEJwr0kHy8xZMDpxbthoBAkPDBps0/view?usp=sharing}{aca}
\end{itemize}

\subsubsection{Semana 08, martes 28/04. Unidad 3. La práctica - Pueblos
originarios}\label{semana-08-martes-2804.-unidad-3.-la-pruxe1ctica---pueblos-originarios}

\begin{enumerate}
\def\labelenumi{\alph{enumi}.}
\setcounter{enumi}{5}
\tightlist
\item
  Pueblos originarios
\end{enumerate}

\begin{itemize}
\item
  Comite DH, NFL,
  \href{https://drive.google.com/file/d/1CP_TpEaI6XXZIb3ACuyqa8XFPuvbdR3f/view?usp=sharing}{aca}
\item
  CSJN, Fallo Coronel, Gustavo Javier,
  \href{https://drive.google.com/file/d/1kslBUCofBjqzYSqUjl-dV5q7sZo_unnV/view?usp=sharing}{aca}
\item
  Inadi, Dictamen Colihueque,
  \href{https://drive.google.com/file/d/1_t9jnOgqHc9qw7PyU6nuJS0_VORXtdLj/view?usp=sharing¨}{aca}
\item
  Video sentencia colihueque
\end{itemize}

\subsubsection{\texorpdfstring{{Semana 08, viernes 01/05. Feriado - Dia
del
Trabajador/a}}{Semana 08, viernes 01/05. Feriado - Dia del Trabajador/a}}\label{semana-08-viernes-0105.-feriado---dia-del-trabajadora}

\subsubsection{Semana 09, martes 05/05 . Unidad 4. Litigio
Constitucional}\label{semana-09-martes-0505-.-unidad-4.-litigio-constitucional}

Estándares de la Corte Suprema para facilitar el acceso a la justicia:

\begin{itemize}
\item
  CSJN
  \href{https://drive.google.com/file/d/1EcDi8VnbuUNH7DlaRPhTMit0IGLLGn0d/view?usp=sharing}{Halabi}
\item
  CSJN
  \href{https://drive.google.com/file/d/186eN_fSToqx4d_W2fRI2pckZA3a5CWCH/view?usp=sharing}{Padec}
\item
  Amparo colectivo
  \href{https://drive.google.com/file/d/1JKvY_Tlo8fV1MUStzKk_WLW_QkCFG1LG/view?usp=sharing}{DADSE}
\end{itemize}

otros ejemplos

\begin{itemize}
\item
  CNCAF, Sala V,
  \href{https://drive.google.com/file/d/1j4N1FxaCWPwVycMh9KS32oLPEqx-vAac/view?usp=sharing}{CELS
  y OTROS s/DNU 70-17}
\item
  CSJN,
  \href{https://drive.google.com/file/d/1kepRDLt6tK6OjOfKD_d4vBkkivAKNYHk/view?usp=sharing}{Quisbert
  Castro}
\end{itemize}

\subsubsection{Semana 09, viernes 08/05. Examen final
presencial}\label{semana-09-viernes-0805.-examen-final-presencial}

\end{document}
